\chapter{Solution Implementation and Evaluation}
\label{sec:implementation}
% ============================================================

% ------------------------------------------------------------
\section{Introduction}
% ------------------------------------------------------------
This chapter presents the implementation of the Sahhty platform, from the adopted technical environment to the development of the different user interfaces. It describes the technologies used, the main implemented functionalities, and the innovative mechanisms integrated into the system, including AI-based pregnancy risk assessment and Wear OS smartwatch integration.

% ------------------------------------------------------------
\section{Environmental Setup}

This section presents all software tools and technologies used in the development of the Sahhty platform, organized by functional category.

% ── Development and Collaboration Tools ──────────────────────
\subsection{Development and Collaboration Tools}
This subsection presents the main tools used during the development and implementation of the Sahhty platform.

\begin{center}
  \includegraphics[height=1.5cm]{chapter4/images/logo_vscode.png}\\[0.3em]
  \small\textbf{Visual Studio Code}
\end{center}

\noindent\textbf{Visual Studio Code} is a lightweight but powerful source code editor developed by Microsoft. It supports a wide range of programming languages and provides built-in tools for debugging, version control, and terminal access. Its extension marketplace allowed us to install plugins for Python, Dart, and Flutter, making it one of our main working environments throughout the project.

\vspace{1em}
\begin{center}
  \includegraphics[height=1.5cm]{chapter4/images/logo_androidstudio.png}\\[0.3em]
  \small\textbf{Android Studio}
\end{center}
\noindent\textbf{Android Studio} is the official integrated development environment for Android application development. It includes a built-in Android emulator that was used extensively to run and test the Flutter mobile application on a virtual device without requiring a physical phone.

\vspace{1em}
\begin{center}
  \includegraphics[height=1.5cm]{chapter4/images/logo_git.png}\\[0.3em]
  \small\textbf{Git}
\end{center}
\noindent\textbf{Git} is a distributed version control system that records every change made to the codebase. It allows developers to work on different features simultaneously, merge their contributions, and revert mistakes at any point. Every commit in our project was tracked using Git, providing a complete and reliable history of the development process.

\vspace{1em}
\begin{center}
  \includegraphics[height=1.5cm]{chapter4/images/logo_github.png}\\[0.3em]
  \small\textbf{GitHub}
\end{center}
\noindent\textbf{GitHub} is an online hosting platform for Git repositories. It served as the central location where the project code is stored and shared between team members. Beyond storage, it also provides tools for tracking issues and reviewing changes through pull requests. The Sahhty source code was hosted on GitHub throughout the entire development period.
\newpage
\vspace{1em}
\begin{center}
  \includegraphics[height=1.5cm]{chapter4/images/logo_postman.png}\\[0.3em]
  \small\textbf{Postman}
\end{center}
\noindent\textbf{Postman} is a tool dedicated to testing HTTP APIs. It provides a graphical interface where the developer can send requests to any endpoint, inspect the response body and headers, and verify that the backend behaves correctly. We used Postman extensively during backend development to validate every API route before integrating it into the Flutter frontend.

\vspace{1em}
\begin{center}
  \includegraphics[height=1.5cm]{chapter4/images/logo_overleaf.png}\\[0.3em]
  \small\textbf{Overleaf}
\end{center}
\noindent\textbf{Overleaf} is a cloud-based \LaTeX{} editor used for writing technical and academic documents. It compiles the document in real time and produces clean, professionally formatted PDFs. We chose it to write this report because it handles formatting, numbering, and bibliographic references automatically.
\vspace{1em}
\begin{center}
  \includegraphics[height=1.5cm]{chapter4/images/Draw.jpeg}\\[0.3em]
  \small\textbf{Overleaf}
\end{center}
\noindent\textbf{Draw.io}(diagrams.net) is a free, open-source diagramming tool used to
design all UML diagrams of the Sahhty platform. It provides a drag-and-drop
interface with built-in UML shape libraries covering use case, class, sequence,
and activity diagrams. Its ability to export diagrams as high-resolution PNG and
SVG files made it straightforward to integrate the resulting diagrams into this
report.
% ── Frontend Technologies ─────────────────────────────────────
\subsection{Frontend Technologies}
This subsection presents the main frontend technologies used for the development of the Sahhty mobile application.
\vspace{0.5em}
\begin{center}
  \includegraphics[height=1.5cm]{chapter4/images/logo_flutter.png}\\[0.3em]
  \small\textbf{Flutter}
\end{center}
\noindent\textbf{Flutter} is an open-source UI toolkit developed by Google for building cross-platform applications from a single codebase using the Dart programming language. The Sahhty mobile application was developed using Flutter 3.x for Android deployment.
\vspace{1em}
\begin{center}
  \includegraphics[height=1.5cm]{chapter4/images/logo_dart.png}\\[0.3em]
  \small\textbf{Dart}
\end{center}
\noindent\textbf{Dart} is a strongly typed, object-oriented programming language developed by Google. It is the language used to write Flutter applications. Its asynchronous programming model, based on \texttt{async/await} and \texttt{Future}, is particularly well-suited for building reactive mobile interfaces that communicate with a remote backend.

\vspace{1em}
\begin{center}
  \includegraphics[height=1.5cm]{chapter4/images/logo_riverpod.png}\\[0.3em]
  \small\textbf{Riverpod}
\end{center}
\noindent\textbf{Riverpod} is a state management library for Flutter. It provides a compile-safe, testable, and reactive approach to managing application state. All data fetching, authentication state, and UI state in the Sahhty frontend are managed using Riverpod providers, ensuring a clean separation between business logic and presentation.

\vspace{1em}
\begin{center}
  \includegraphics[height=1.5cm]{chapter4/images/logo_dio.png}\\[0.3em]
  \small\textbf{Dio}
\end{center}
\noindent\textbf{Dio} is a powerful HTTP client for Dart. It simplifies sending requests to a REST API and handling responses, with built-in support for interceptors, request cancellation, and file uploads. We configured Dio with a base URL and an interceptor that automatically attaches the JWT access token to every outgoing authenticated request.

\vspace{1em}
\begin{center}
  \includegraphics[height=1.5cm]{chapter4/images/logo_gorouter.png}\\[0.3em]
  \small\textbf{Go Router}
\end{center}
\noindent\textbf{Go Router} is the official routing package for Flutter applications. It provides declarative, URL-based navigation with support for route guards and nested routes. We used it to implement role-based navigation, redirecting patients and doctors to their respective home screens after login.

\vspace{1em}
\begin{center}
\includegraphics[height=1.5cm]{chapter4/images/logo_flutter.png}\\[0.3em]
  \small\textbf{Flutter Secure Storage}
\end{center}

\noindent\textbf{Flutter Secure Storage} is a package that stores sensitive data, such as JWT tokens, in the platform's secure keystore (Android Keystore on Android). This ensures that authentication credentials persist across app restarts without being exposed to other applications.

\vspace{1em}
\begin{center}
  \includegraphics[height=1.5cm]{chapter4/images/logo_firebase.png}\\[0.3em]
  \small\textbf{Firebase Messaging}
\end{center}
\noindent\textbf{Firebase Messaging} is the Flutter integration of Firebase Cloud Messaging (FCM). It is used to receive push notifications on the Android device when the application is running in the background or fully closed, ensuring that critical health alerts and appointment reminders always reach the patient.

\vspace{1em}
\begin{center}

  \small\textbf{Flutter Animate}
\end{center}
\noindent\textbf{Flutter Animate} is an animation library that provides a declarative API for adding smooth transitions and micro-animations to Flutter widgets. We used it throughout the application to deliver a polished and engaging user experience.

\vspace{1em}
\begin{center}
  \includegraphics[height=1.5cm]{chapter4/images/logo_googlefonts.png}\\[0.3em]
  \small\textbf{Google Fonts}
\end{center}
\noindent\textbf{Google Fonts} and \textbf{Iconsax} are packages used for typography and iconography. Google Fonts provides access to the Poppins typeface used consistently across the application. Iconsax provides a comprehensive set of clean, modern icons that replace system icons for a cohesive visual identity.

\vspace{1em}
\begin{center}
  \includegraphics[height=1.5cm]{chapter4/images/logo_fluttermap.png}\\[0.3em]
  \small\textbf{Flutter Map}
\end{center}
\noindent\textbf{Flutter Map} and \textbf{Latlong2} are packages used to render interactive maps powered by OpenStreetMap tiles. They are used in the doctor discovery screen to display the geographic locations of doctors on an interactive map, enabling patients to find nearby specialists.

% ── Backend Technologies ──────────────────────────────────────
\subsection{Backend Technologies}
This subsection presents the main backend technologies used for the development of the Sahhty platform server and API infrastructure.
\vspace{0.5em}
\begin{center}
  \includegraphics[height=1.5cm]{chapter4/images/logo_python.png}\\[0.3em]
  \small\textbf{Python}
\end{center}
\noindent\textbf{Python} is a high-level, general-purpose programming language valued for its readability and large ecosystem of libraries. It is widely adopted in web development, data science, and artificial intelligence. We used Python~3.13 as the language for the entire backend of Sahhty.

\vspace{1em}
\begin{center}
  \includegraphics[height=1.5cm]{chapter4/images/logo_django.png}\\[0.3em]
  \small\textbf{Django}
\end{center}
\noindent\textbf{Django} is a high-level Python web framework that encourages rapid development and clean, pragmatic design. It provides a built-in ORM, admin panel, and authentication system. The entire Sahhty backend is structured around Django~5.x with multiple domain-specific applications covering users, patients, doctors, appointments, measurements, alerts, medications, and medical files.

\vspace{1em}
\begin{center}
  \includegraphics[height=1.5cm]{chapter4/images/logo_drf.png}\\[0.3em]
  \small\textbf{Django REST Framework}
\end{center}
\noindent\textbf{Django REST Framework (DRF)} is a toolkit for building RESTful APIs with Django. It provides serializers, viewsets, and authentication mechanisms used to implement the Sahhty backend API endpoints.
\vspace{1em}
\begin{center}
  \includegraphics[height=1.5cm]{chapter4/images/logo_simplejwt.png}\\[0.3em]
  \small\textbf{Simple JWT}
\end{center}
\noindent\textbf{Simple JWT} is a Django REST Framework extension that provides JSON Web Token authentication. It issues access and refresh tokens upon login, and validates token signatures on every protected request. The Sahhty backend uses short-lived access tokens with a refresh mechanism to balance security and usability.

\vspace{1em}
\begin{center}
  \includegraphics[height=1.5cm]{chapter4/images/logo_apscheduler.png}\\[0.3em]
  \small\textbf{APScheduler}
\end{center}
\noindent\textbf{APScheduler} is a Python library for scheduling background jobs. It is used in the Sahhty backend to run periodic tasks such as sending medication reminder notifications, appointment reminder alerts, and pregnancy monitoring checks at defined time intervals.

\vspace{1em}
\begin{center}
  \includegraphics[height=1.5cm]{chapter4/images/logo_websocket.png}\\[0.3em]
  \small\textbf{WebSocket}
\end{center}
\noindent\textbf{WebSocket} is a communication protocol that provides a persistent
full-duplex channel over a single TCP connection, allowing the server to push data
to the client instantly without waiting for a request. Unlike HTTP, WebSocket keeps
the connection open after the initial handshake, eliminating the need for polling
and enabling real-time communication. In the Sahhty platform, WebSocket is used to
deliver instant notifications for appointment events and medical file access
requests, with each connection authenticated via a JWT token passed during the
handshake.

\vspace{1em}
\begin{center}
  \includegraphics[height=1.5cm]{chapter4/images/logo_channels.png}\\[0.3em]
  \small\textbf{Django Channels}
\end{center}
\noindent\textbf{Django Channels} extends Django to handle asynchronous protocols
beyond standard HTTP, most notably WebSockets. It is used in the Sahhty backend
to maintain persistent bidirectional connections with connected clients, enabling
real-time delivery of notifications for appointment events and medical file access
requests without requiring the client to poll the server.

\vspace{1em}
\begin{center}
  \includegraphics[height=1.5cm]{chapter4/images/logo_memurai.png}\\[0.3em]
  \small\textbf{Memurai}
\end{center}
\noindent\textbf{Memurai} is a native Redis-compatible in-memory data store for
Windows. It serves as the channel layer backend required by Django Channels to
route WebSocket messages between the server and connected clients. Every real-time
notification dispatched by the backend passes through Memurai before being
delivered to the target WebSocket connection.

\vspace{1em}
\begin{center}
  \includegraphics[height=1.5cm]{chapter4/images/logo_sklearn.png}\\[0.3em]
  \small\textbf{Scikit-learn}
\end{center}
\noindent\textbf{Scikit-learn} is a Python machine learning library used alongside XGBoost for data preprocessing, feature engineering, and model evaluation during the training phase of the pregnancy risk assessment model.

% ── Database and Infrastructure ───────────────────────────────
\subsection{Database and Infrastructure}
This subsection presents the database technologies and infrastructure services used to support the deployment and execution of the Sahhty platform.
\vspace{0.5em}
\begin{center}
  \includegraphics[height=1.5cm]{chapter4/images/logo_postgresql.png}\\[0.3em]
  \small\textbf{PostgreSQL}
\end{center}
\noindent\textbf{PostgreSQL} is an advanced open-source relational database system. It serves as the primary data store for all Sahhty platform data, including user accounts, patient profiles, medical records, appointments, and health measurements. Its support for transactions, foreign key constraints, and complex queries ensures data integrity across all operations.

\vspace{1em}
\begin{center}
  \includegraphics[height=1.5cm]{chapter4/images/logo_neon.png}\\[0.3em]
  \small\textbf{Neon PostgreSQL}
\end{center}
\noindent\textbf{Neon PostgreSQL} is a serverless PostgreSQL hosting service used to host the production database instance. It provides automatic scaling, branching, and connection pooling, making it well-suited for a growing mobile health application.

% ── External Services ─────────────────────────────────────────
\subsection{External Services}
This subsection presents the external services integrated into the Sahhty platform to support notifications, communication, and additional system functionalities.
\vspace{0.5em}
\begin{center}
  \includegraphics[height=1.5cm]{chapter4/images/logo_firebase.png}\\[0.3em]
  \small\textbf{Firebase Cloud Messaging}
\end{center}
\noindent\textbf{Firebase Cloud Messaging (FCM)} is a push notification service used to deliver real-time alerts from the Django backend to Android devices.
\vspace{1em}
\begin{center}
  \includegraphics[height=1.5cm]{chapter4/images/logo_brevo.png}\\[0.3em]
  \small\textbf{Brevo}
\end{center}
\noindent\textbf{Brevo} is a transactional email service used to deliver account verification OTPs and system notifications by email. It provides a reliable SMTP relay that handles email delivery, tracking, and bounce management.

\vspace{1em}
\begin{center}
  \includegraphics[height=1.5cm]{chapter4/images/logo_openstreetmap.png}\\[0.3em]
  \small\textbf{OpenStreetMap}
\end{center}
\noindent\textbf{OpenStreetMap / Tile Servers} provide the map tile layer rendered in the Flutter map widget. The map data is publicly available and free of charge, making it the appropriate choice for an academic prototype.

\vspace{1em}
\begin{center}
  \includegraphics[height=1.5cm]{chapter4/images/logo_groq.png}\\[0.3em]
  \small\textbf{Groq}
\end{center}
\noindent\textbf{Groq} is a high-performance AI inference platform that provides ultra-fast access to large language models (LLMs) through a simple REST API. In the Sahhty platform, Groq is used to power natural language processing features that require near real-time responses, leveraging its LPU (Language Processing Unit) hardware to deliver inference speeds significantly faster than conventional GPU-based services. Its free tier made it an appropriate choice for an academic prototype requiring intelligent text generation capabilities.
\newpage
\section{Interface Presentation}
This section presents the main interfaces of the Sahhty platform and the different functionalities available for patients and doctors.

\subsection{Logo Design}
\begin{figure}[H]
  \centering
  \includegraphics[width=0.25\textwidth]{chapter4/images/logo.png}
  \caption{Sahhty application logo}
  \label{fig:logo}
\end{figure}
The Sahhty app logo features a heart shape with a smooth gradient blending turquoise green and blue, symbolizing health, vitality, and care. At the center of the heart, a white ECG/heartbeat line represents medical monitoring and life. The typography uses the name "Sahhty" in a modern, rounded dark font, conveying trust and professionalism. The overall design is clean, minimalist, and perfectly suited for a health and maternity mobile application.


\subsection{Color Design}

The visual identity of the Sahhty platform was designed to balance medical professionalism, emotional comfort, and usability. The adopted color palette helps create a reassuring experience for pregnant women while ensuring clear distinction between alerts, medical states, and user roles. The design rationale of the adopted color system is presented in Table~\ref{tab:colors}.
\clearpage
\renewcommand{\arraystretch}{1.5}

\begin{longtable}{|p{2.5cm}|p{2cm}|p{3.5cm}|p{5.5cm}|}
\caption{Sahhty application color palette and design rationale}
\label{tab:colors} \\
\hline
\rowcolor{pink!15}
\textbf{Color Name} & \textbf{Preview} & \textbf{Main Usage} & \textbf{Design Rationale} \\
\hline
\endfirsthead
\hline
\rowcolor{pink!15}
\textbf{Color Name} & \textbf{Preview} & \textbf{Main Usage} & \textbf{Design Rationale} \\
\hline
\endhead
Rose / Pink &
\cellcolor[HTML]{E91E8C} &
Primary brand identity and women patient interface &
Conveys maternity, warmth, and emotional comfort while strengthening the platform identity. \\
\hline
Dark Rose &
\cellcolor[HTML]{C2185B} &
Buttons, pressed states, and gradients &
Provides visual depth and interactive feedback while preserving interface consistency. \\
\hline
Light Green &
\cellcolor[HTML]{A5D6A7} &
Doctor interface and confirmation states &
Symbolizes health, validation, and trust, while visually separating doctor-related interfaces. \\
\hline
Dark Blue &
\cellcolor[HTML]{0D47A1} &
Primary brand identity and men patient interface &
Conveys trust, professionalism and reliability, creating a strong and serious medical identity for male patients. \\
\hline
Orange / Amber &
\cellcolor[HTML]{FFB300} &
Warning dialogs and medium-risk alerts &
Represents caution and draws user attention without creating panic. \\
\hline
Red &
\cellcolor[HTML]{E53935} &
Critical alerts and error states &
Indicates urgency and dangerous medical conditions requiring immediate attention. \\
\hline
White &
\cellcolor[HTML]{FFFFFF} &
Backgrounds and content areas &
Improves readability and creates a clean medical appearance. \\
\hline
Dark Grey / Black &
\cellcolor[HTML]{424242} &
Text and medical labels &
Ensures high contrast and accessibility for healthcare information display. \\
\hline
Blue &
\cellcolor[HTML]{42A5F5} &
Blood glucose measurement cards &
Associated with precision and clinical measurements. \\
\hline
Purple &
\cellcolor[HTML]{8E24AA} &
Heart rate measurement cards &
Differentiates cardiovascular indicators from other vital measurements. \\
\hline
Gradient Overlay &
\cellcolor[HTML]{D1C4E9} &
dashboard backgrounds &
Enhances visual aesthetics while maintaining text readability over images. \\
\hline
% ── LOGO COLORS ──────────────────────────
Logo Turquoise &
\cellcolor[HTML]{4DB6AC} &
Logo heart and brand identity &
Represents freshness, care and vitality at the core of the Sahhty brand. \\
\hline
Logo Blue &
\cellcolor[HTML]{1E88E5} &
Logo heart and brand identity &
Conveys trust, professionalism and medical reliability in the brand identity. \\
\hline
\end{longtable}
% ── Interface Overview ────────────────────────────────────────
\subsection{Interface Overview}

This section presents the main interfaces of the Sahhty application organized by functional area.
% ── Authentication ────────────────────────────────────────────
\subsection{Authentication}
This section presents the authentication interfaces of the Sahhty platform, including language selection, patient and doctor registration, OTP verification, login, and two-factor authentication.
\subsubsection{Splash Screen and Language Selection}

At first launch, the user is presented with a language selection screen supporting both French and English. The selected language is then applied globally across the application. Figures~\ref{fig:splash}a and~\ref{fig:splash}b illustrate the splash screen and the language selection interface.

\begin{figure}[H]
  \centering
  \begin{subfigure}[b]{0.30\textwidth}
    \includegraphics[width=\textwidth]{chapter4/images/splash_screen.jpeg}
    \caption{Splash screen}
  \end{subfigure}
\hspace{0.02\textwidth}
  \begin{subfigure}[b]{0.30\textwidth}
    \includegraphics[width=\textwidth]{chapter4/images/language_selection.jpeg}
    \caption{Language selection}
  \end{subfigure}
  \caption{Application startup interfaces}
  \label{fig:splash}
\end{figure}

\subsubsection{Patient Registration}
The patient registration process consists of two steps covering personal, medical, and pregnancy-related information. Figures~\ref{fig:patient_registration}a--e illustrate the registration interfaces.

\begin{figure}[H]
\centering

\begin{subfigure}[b]{0.30\textwidth}
    \includegraphics[width=\textwidth]{chapter4/images/register_step1.jpeg}
    \caption{Personal information}
\end{subfigure}
\hfill
\begin{subfigure}[b]{0.30\textwidth}
    \includegraphics[width=\textwidth]{chapter4/images/register_step1_1.jpeg}
    \caption{Account details}
\end{subfigure}
\hfill
\begin{subfigure}[b]{0.30\textwidth}
    \includegraphics[width=\textwidth]{chapter4/images/register_step1_2.jpeg}
    \caption{OTP verification}
\end{subfigure}
\hfill
\begin{subfigure}[b]{0.30\textwidth}
    \includegraphics[width=\textwidth]{chapter4/images/register_step2.jpeg}
    \caption{Medical information}
\end{subfigure}
\hspace{0.02\textwidth}
\begin{subfigure}[b]{0.30\textwidth}
    \includegraphics[width=\textwidth]{chapter4/images/register_step3.jpeg}
    \caption{Pregnancy information}
\end{subfigure}

\caption{Patient registration workflow}
\label{fig:patient_registration}
\end{figure}

\subsubsection{Doctor Registration}

The doctor registration process allows healthcare professionals to create and configure their professional profile by providing speciality, address, experience, consultation fees, and biography information. After registration, the account remains pending until administrator verification. Figures~\ref{fig:reg_doc}a--c illustrate the doctor registration interfaces.

\begin{figure}[H]
\centering

\begin{subfigure}[b]{0.3\textwidth}
    \includegraphics[width=\textwidth]{chapter4/images/register_doctor.jpeg}
    \caption{Professional information}
\end{subfigure}
\hfill
\begin{subfigure}[b]{0.3\textwidth}
    \includegraphics[width=\textwidth]{chapter4/images/register_doctor_1.jpeg}
    \caption{Profile configuration}
\end{subfigure}
\hfill
\begin{subfigure}[b]{0.3\textwidth}
    \includegraphics[width=\textwidth]{chapter4/images/register_doctor_2.jpeg}
    \caption{Verification pending}
\end{subfigure}

\caption{Doctor registration workflow}
\label{fig:reg_doc}
\end{figure}

\subsubsection{OTP Verification}

After registration, the user receives a six-digit OTP code by email to verify the account before activation. Figures~\ref{fig:otp}a and~\ref{fig:otp}b illustrate the OTP verification interface.

\begin{figure}[H]
\centering

\begin{subfigure}[b]{0.30\textwidth}
    \includegraphics[width=\textwidth]{chapter4/images/otp_screen.jpeg}
    \caption{OTP input screen}
\end{subfigure}
\hspace{0.02\textwidth}
\begin{subfigure}[b]{0.30\textwidth}
    \includegraphics[width=\textwidth]{chapter4/images/otp_screen_1.jpeg}
    \caption{Successful verification}
\end{subfigure}

\caption{Email OTP verification}
\label{fig:otp}
\end{figure}

\subsubsection{Login and Two-Factor Authentication}

The login interface allows users to authenticate using their email and password. If two-factor authentication (2FA) is enabled, an additional verification step is required before accessing the application. Figures~\ref{fig:2fa}a and~\ref{fig:2fa}b illustrate the authentication interfaces.

\begin{figure}[H]
\centering

\begin{subfigure}[b]{0.30\textwidth}
    \includegraphics[width=\textwidth]{chapter4/images/twofa_screen.jpeg}
    \caption{Login screen}
\end{subfigure}
\hspace{0.02\textwidth}
\begin{subfigure}[b]{0.30\textwidth}
    \includegraphics[width=\textwidth]{chapter4/images/twofa_screen_1.jpeg}
    \caption{2FA verification}
\end{subfigure}

\caption{User authentication interfaces}
\label{fig:2fa}
\end{figure}

\subsection{Patient Home Screen and Pregnancy Dashboard}

The patient home screen serves as the main dashboard of the application by presenting pregnancy progression, recent health measurements, upcoming appointments, and recent alerts in a single interface. Figures~\ref{fig:home}a--c illustrate the different sections of the patient dashboard.

\begin{figure}[H]
\centering

\begin{subfigure}[b]{0.30\textwidth}
    \includegraphics[width=\textwidth]{chapter4/images/home_patient.jpeg}
    \caption{Pregnancy dashboard}
\end{subfigure}
\hfill
\begin{subfigure}[b]{0.30\textwidth}
    \includegraphics[width=\textwidth]{chapter4/images/home_patient_1.jpeg}
    \caption{Health measurements}
\end{subfigure}
\hfill
\begin{subfigure}[b]{0.30\textwidth}
    \includegraphics[width=\textwidth]{chapter4/images/home_patient_2.jpeg}
    \caption{Appointments and alerts}
\end{subfigure}

\caption{Patient home screen}
\label{fig:home}
\end{figure}
\subsubsection{Pregnancy Tracking Screen}
The pregnancy tracking screen provides detailed information about the current pregnancy, including the estimated due date, pregnancy duration, and visual progression indicators. A distinctive feature of the pregnancy tracking screen is the fruit-based fetal size visualization, mapping each stage of pregnancy to a familiar analogy: a seed (weeks 1--4), blueberry (weeks 5--8), grape (weeks 9--12), lemon (weeks 13--16), banana (weeks 17--20), mango (weeks 21--24), coconut (weeks 25--28), melon (weeks 29--32), pumpkin (weeks 33--36), and watermelon (week 37+). The corresponding emoji is displayed prominently in the hero card with a subtle pulsing animation, giving the patient a clear and emotionally engaging sense of her baby's development at every stage. Figures~\ref{fig:preg}a and~\ref{fig:preg}b illustrate the pregnancy tracking interfaces.\ref{fig:preg}c illustrate the Fetal size visualization by week .

\begin{figure}[H]
\centering
\begin{subfigure}[b]{0.30\textwidth}
    \includegraphics[width=\textwidth]{chapter4/images/pregnancy_screen.jpeg}
    \caption{Pregnancy overview}
\end{subfigure}
\hspace{0.02\textwidth}
\begin{subfigure}[b]{0.30\textwidth}
    \includegraphics[width=\textwidth]{chapter4/images/pregnancy_screen_1.jpeg}
    \caption{Pregnancy progression}
\end{subfigure}

\vspace{1em}

\begin{subfigure}[b]{0.75\textwidth}
    \centering
    \includegraphics[width=\textwidth]{chapter4/images/fruit_timeline.png}
    \caption{Fetal size visualization by week}
\end{subfigure}

\caption{Pregnancy tracking screen and fetal size visualization}
\label{fig:preg}
\end{figure}

\subsection{Health Measurements}

This section presents the interfaces related to patient health measurements, including measurement history management, manual measurement submission, AI-based risk alerts, and smartwatch integration for automatic heart rate synchronization.Figure~\ref{fig:meas_list} illustrates the measurements history interface.

\subsubsection{Measurements List}
The measurements screen displays the patient's recorded health measurements with filtering, sorting, and pagination functionalities. Figure~\ref{fig:meas_list} illustrates the measurements history interface.

\begin{figure}[H]
  \centering
  \begin{subfigure}[b]{0.30\textwidth}
    \includegraphics[width=\textwidth]{chapter4/images/measurements_list.jpeg}
    \caption{Measurements history}
  \end{subfigure}

  \caption{Health measurements list}
  \label{fig:meas_list}
\end{figure}

\subsubsection{Add Measurement}

The add measurement interface allows the patient to record new health readings and automatically triggers the AI-based pregnancy risk assessment process. Figures~\ref{fig:add_meas}a and~\ref{fig:add_meas}b illustrate the measurement submission form and the generated high-risk alert.

\begin{figure}[H]
  \centering

  \begin{subfigure}[b]{0.30\textwidth}
    \includegraphics[width=\textwidth]{chapter4/images/add_measurement.jpeg}
    \caption{Measurement form}
  \end{subfigure}
\hspace{0.02\textwidth}
  \begin{subfigure}[b]{0.30\textwidth}
    \includegraphics[width=\textwidth]{chapter4/images/risk_alert_high.jpeg}
    \caption{High-risk alert}
  \end{subfigure}

  \caption{Health measurement and AI-generated risk alert}
  \label{fig:add_meas}
\end{figure}

\subsubsection{Smartwatch Integration}

The Sahhty platform integrates with Wear OS smartwatches to automatically collect and synchronize heart rate measurements with the mobile application. Figures~\ref{fig:watch}a and~\ref{fig:watch}b illustrate the smartwatch integration interfaces.

\begin{figure}[H]
  \centering
  \begin{subfigure}[b]{0.30\textwidth}
    \includegraphics[width=\textwidth]{chapter4/images/smartwatch_screen.jpeg}
    \caption{Smartwatch synchronization}
  \end{subfigure}
\hspace{0.02\textwidth}
  \begin{subfigure}[b]{0.30\textwidth}
    \includegraphics[width=\textwidth]{chapter4/images/smartwatch_screen_1.jpeg}
    \caption{Heart rate monitoring}
  \end{subfigure}

  \caption{Wear OS smartwatch integration}
  \label{fig:watch}
\end{figure}

\subsection{Notifications and Alerts}

The notifications interface displays health and system alerts generated by the backend, with filtering by category and severity level. Figures~\ref{fig:alerts}a and~\ref{fig:alerts}b illustrate the patient alert management interfaces.

\begin{figure}[H]
  \centering
  \begin{subfigure}[b]{0.30\textwidth}
    \includegraphics[width=\textwidth]{chapter4/images/alerts_list.jpeg}
    \caption{Alerts list}
  \end{subfigure}
\hspace{0.02\textwidth}
  \begin{subfigure}[b]{0.30\textwidth}
    \includegraphics[width=\textwidth]{chapter4/images/alerts_list_1.jpeg}
    \caption{Alert details}
  \end{subfigure}

  \caption{Patient notifications and alerts}
  \label{fig:alerts}
\end{figure}


\subsection{Doctor Discovery and Appointment Booking}

This section presents the interfaces related to doctor search, speciality filtering, geographic visualization, and appointment booking.

\subsubsection{Find a Doctor}

The doctor discovery interface allows patients to search for doctors by name or speciality using both list and map views. Figure~\ref{fig:doctor_discovery}a presents the searchable doctor list with speciality filters, while Figure~\ref{fig:doctor_discovery}b displays the geographic map view of available doctors.

\begin{figure}[H]
  \centering

  \begin{subfigure}[b]{0.30\textwidth}
    \includegraphics[width=\textwidth]{chapter4/images/find_doctor_list.jpeg}
    \caption{Doctor list view}
  \end{subfigure}
\hspace{0.02\textwidth}
  \begin{subfigure}[b]{0.30\textwidth}
    \includegraphics[width=\textwidth]{chapter4/images/find_doctor_map.jpeg}
    \caption{Doctor map view}
  \end{subfigure}

  \caption{Doctor discovery interfaces}
  \label{fig:doctor_discovery}
\end{figure}

\subsubsection{Doctor Profile and Appointment Booking}

The doctor profile interface presents detailed professional information, consultation schedules, location, and appointment availability. Patients can select an available date and confirm an appointment directly from the interface. Figures~\ref{fig:doc_profile}a--c illustrate the doctor profile and appointment booking workflow.

\begin{figure}[H]
  \centering

  \begin{subfigure}[b]{0.30\textwidth}
    \includegraphics[width=\textwidth]{chapter4/images/doctor_profile.jpeg}
    \caption{Doctor profile}
  \end{subfigure}
  \hfill
  \begin{subfigure}[b]{0.30\textwidth}
    \includegraphics[width=\textwidth]{chapter4/images/doctor_profile_1.jpeg}
    \caption{Schedule and availability}
  \end{subfigure}
  \hfill
  \begin{subfigure}[b]{0.30\textwidth}
    \includegraphics[width=\textwidth]{chapter4/images/doctor_profile_2.jpeg}
    \caption{Appointment booking}
  \end{subfigure}

  \caption{Doctor profile and appointment booking}
  \label{fig:doc_profile}
\end{figure}

\subsubsection{My Appointments}

The appointments interface displays all upcoming and previous appointments of the patient, including doctor information, appointment date, and current status. Figure~\ref{fig:appts} illustrates the patient appointment management interface.

\begin{figure}[H]
  \centering

  \begin{subfigure}[b]{0.30\textwidth}
    \includegraphics[width=\textwidth]{chapter4/images/appointments_list.jpeg}
    \caption{Appointments list}
  \end{subfigure}

  \caption{Patient appointments}
  \label{fig:appts}
\end{figure}

% ── Medications & DCI ─────────────────────────────────────────
\subsection{Medications and Drug Safety Analysis}

This section presents the interfaces related to treatment management, medication search, pregnancy drug safety classification, and DCI interaction analysis.

\subsubsection{My Treatments}

The treatments interface displays the patient's active medications, dosage schedules, treatment duration, and allergy incompatibility warnings. Figures~\ref{fig:meds}a--c illustrate the patient treatment management interfaces.

\begin{figure}[H]
  \centering

  \begin{subfigure}[b]{0.30\textwidth}
    \includegraphics[width=\textwidth]{chapter4/images/medications_list.jpeg}
    \caption{Treatment list}
  \end{subfigure}
  \hfill
  \begin{subfigure}[b]{0.30\textwidth}
    \includegraphics[width=\textwidth]{chapter4/images/medications_list_1.jpeg}
    \caption{Treatment details}
  \end{subfigure}
  \hfill
  \begin{subfigure}[b]{0.30\textwidth}
    \includegraphics[width=\textwidth]{chapter4/images/medications_list_2.jpeg}
    \caption{Allergy warning}
  \end{subfigure}

  \caption{Patient treatment management}
  \label{fig:meds}
\end{figure}

\subsubsection{Medication Search and Safety Analysis}

The medication search functionality allows users to search the national medication database and access pregnancy risk classification, trimester-specific information, and DCI interaction analysis. Figures~\ref{fig:med_search} and~\ref{fig:med_detail} illustrate the medication search and detailed safety analysis interfaces.

\begin{figure}[H]
  \centering

  \begin{subfigure}[b]{0.30\textwidth}
    \includegraphics[width=\textwidth]{chapter4/images/medication_search.jpeg}
    \caption{Medication search}
  \end{subfigure}
  \hspace{0.01\textwidth}
  \begin{subfigure}[b]{0.30\textwidth}
    \includegraphics[width=\textwidth]{chapter4/images/medication_search_1.jpeg}
    \caption{Search results}
  \end{subfigure}
  \hspace{0.01\textwidth}
  \begin{subfigure}[b]{0.30\textwidth}
    \includegraphics[width=\textwidth]{chapter4/images/medication_search_2.jpeg}
    \caption{Pregnancy classification}
  \end{subfigure}

  \vspace{0.5cm}

  \begin{subfigure}[b]{0.30\textwidth}
    \includegraphics[width=\textwidth]{chapter4/images/medication_detail.jpeg}
    \caption{Medication details}
  \end{subfigure}
  \hspace{0.01\textwidth}
  \begin{subfigure}[b]{0.30\textwidth}
    \includegraphics[width=\textwidth]{chapter4/images/medication_detail_1.jpeg}
    \caption{Trimester risks}
  \end{subfigure}
  \hspace{0.01\textwidth}
  \begin{subfigure}[b]{0.30\textwidth}
    \includegraphics[width=\textwidth]{chapter4/images/medication_detail_2.jpeg}
    \caption{DCI interactions}
  \end{subfigure}

  \caption{Medication search and safety analysis interfaces}
  \label{fig:medication_analysis}
\end{figure}

\subsection{Medical File Management}

This section presents the interfaces related to medical file storage, document visualization, and doctor access management.

\subsubsection{Patient Medical Files}

The medical files interface allows patients to upload, organize, and consult their personal medical documents by category. Figure~\ref{fig:med_files} illustrates the medical file management interface.

\begin{figure}[H]
  \centering

  \begin{subfigure}[b]{0.30\textwidth}
    \includegraphics[width=\textwidth]{chapter4/images/medical_files.jpeg}
    \caption{Medical files overview}
  \end{subfigure}

  \caption{Patient medical files}
  \label{fig:med_files}
\end{figure}

\subsubsection{Medical File Access Management}

The access management interface allows patients to control which doctors are authorized to access their medical files. Figures~\ref{fig:access}a--c illustrate the authorization management interfaces.

\begin{figure}[H]
  \centering

  \begin{subfigure}[b]{0.30\textwidth}
    \includegraphics[width=\textwidth]{chapter4/images/medical_files_access.jpeg}
    \caption{Authorized doctors}
  \end{subfigure}
  \hfill
  \begin{subfigure}[b]{0.30\textwidth}
    \includegraphics[width=\textwidth]{chapter4/images/medical_files_access_1.jpeg}
    \caption{Pending requests}
  \end{subfigure}
  \hfill
  \begin{subfigure}[b]{0.30\textwidth}
    \includegraphics[width=\textwidth]{chapter4/images/medical_files_access_2.jpeg}
    \caption{Access management}
  \end{subfigure}

  \caption{Medical file access management}
  \label{fig:access}
\end{figure}

\subsection{Doctor Dashboard}

This section presents the main interfaces available to doctors for managing appointments, patient information, alerts, and professional activity.

\subsubsection{Doctor Home Screen}

The doctor home screen provides a dashboard summarizing appointments, patient activity, alerts, and schedule management features. Figures~\ref{fig:home_doc}a--c illustrate the main doctor dashboard interfaces.

\begin{figure}[H]
  \centering

  \begin{subfigure}[b]{0.30\textwidth}
    \includegraphics[width=\textwidth]{chapter4/images/home_doctor.jpeg}
    \caption{Doctor dashboard}
  \end{subfigure}
  \hfill
  \begin{subfigure}[b]{0.30\textwidth}
    \includegraphics[width=\textwidth]{chapter4/images/home_doctor_1.jpeg}
    \caption{Statistics overview}
  \end{subfigure}
  \hfill
  \begin{subfigure}[b]{0.30\textwidth}
    \includegraphics[width=\textwidth]{chapter4/images/home_doctor_2.jpeg}
    \caption{Quick actions}
  \end{subfigure}

  \caption{Doctor home screen}
  \label{fig:home_doc}
\end{figure}
\subsubsection{Doctor Appointments}

The doctor appointments interface allows doctors to manage incoming and previous appointments, including appointment confirmation and cancellation actions. Figures~\ref{fig:doc_appts}a and~\ref{fig:doc_appts}b illustrate the appointment management interfaces.

\begin{figure}[H]
  \centering

  \begin{subfigure}[b]{0.30\textwidth}
    \includegraphics[width=\textwidth]{chapter4/images/doctor_appointments.jpeg}
    \caption{Appointments overview}
  \end{subfigure}
\hspace{0.02\textwidth}
  \begin{subfigure}[b]{0.30\textwidth}
    \includegraphics[width=\textwidth]{chapter4/images/doctor_appointments_1.jpeg}
    \caption{Appointment management}
  \end{subfigure}

  \caption{Doctor appointment management}
  \label{fig:doc_appts}
\end{figure}

\subsubsection{Patient Medical File Access}

The medical file access interface allows doctors to request and manage access to patient medical records, including documents, measurements, and treatments. Figures~\ref{fig:doc_access}a--c illustrate the medical file access workflow.

\begin{figure}[H]
  \centering

  \begin{subfigure}[b]{0.30\textwidth}
    \includegraphics[width=\textwidth]{chapter4/images/doctor_medical_access.jpeg}
    \caption{Patient search}
  \end{subfigure}
  \hfill
  \begin{subfigure}[b]{0.30\textwidth}
    \includegraphics[width=\textwidth]{chapter4/images/doctor_medical_access_1.jpeg}
    \caption{Access request}
  \end{subfigure}
  \hfill
  \begin{subfigure}[b]{0.30\textwidth}
    \includegraphics[width=\textwidth]{chapter4/images/doctor_medical_access_2.jpeg}
    \caption{Medical file access}
  \end{subfigure}

  \caption{Doctor medical file access management}
  \label{fig:doc_access}
\end{figure}
\subsubsection{Doctor Alerts}

The doctor alerts interface centralizes appointment notifications, patient access requests, and system-related events through categorized tabs. Figures~\ref{fig:doc_alerts}a--c illustrate the doctor notification management interfaces.

\begin{figure}[H]
  \centering

  \begin{subfigure}[b]{0.30\textwidth}
    \includegraphics[width=\textwidth]{chapter4/images/doctor_alerts.png}
    \caption{Appointment alerts}
  \end{subfigure}
  \hfill
  \begin{subfigure}[b]{0.30\textwidth}
    \includegraphics[width=\textwidth]{chapter4/images/doctor_alerts_1.png}
    \caption{Patient requests}
  \end{subfigure}
  \hfill
  \begin{subfigure}[b]{0.30\textwidth}
    \includegraphics[width=\textwidth]{chapter4/images/doctor_alerts_2.png}
    \caption{System notifications}
  \end{subfigure}

  \caption{Doctor alerts and notifications}
  \label{fig:doc_alerts}
\end{figure}

% ── Doctor Schedule Management ────────────────────────────────
\subsection{Doctor Schedule Management}

The schedule management interface allows doctors to configure their weekly availability, consultation hours, leave days, and appointment capacity. Figures~\ref{fig:schedule}a--c illustrate the doctor scheduling interfaces.

\begin{figure}[H]
  \centering

  \begin{subfigure}[b]{0.30\textwidth}
    \includegraphics[width=\textwidth]{chapter4/images/doctor_schedule.jpeg}
    \caption{Weekly schedule}
  \end{subfigure}
  \hfill
  \begin{subfigure}[b]{0.30\textwidth}
    \includegraphics[width=\textwidth]{chapter4/images/doctor_schedule_1.jpeg}
    \caption{Availability configuration}
  \end{subfigure}
  \hfill
  \begin{subfigure}[b]{0.30\textwidth}
    \includegraphics[width=\textwidth]{chapter4/images/doctor_schedule_2.jpeg}
    \caption{Leave management}
  \end{subfigure}

  \caption{Doctor schedule management}
  \label{fig:schedule}
\end{figure}
% ── Profile and Settings ──────────────────────────────────────
\subsection{Profile and Settings}

This section presents the interfaces related to patient profile management and application settings.

\subsubsection{Patient Profile}

The patient profile interface displays personal information, medical data, and menstrual cycle details organized into expandable sections. Figure~\ref{fig:profile} illustrates the patient profile screen.

\begin{figure}[H]
  \centering

  \begin{subfigure}[b]{0.30\textwidth}
    \includegraphics[width=\textwidth]{chapter4/images/patient_profile.jpeg}
    \caption{Patient profile}
  \end{subfigure}

  \caption{Patient profile interface}
  \label{fig:profile}
\end{figure}
\subsubsection{Settings}

The settings interface provides access to language selection, notification preferences, two-factor authentication, and account management features. Figures~\ref{fig:settings}a--c illustrate the application settings interfaces.

\begin{figure}[H]
  \centering

  \begin{subfigure}[b]{0.30\textwidth}
    \includegraphics[width=\textwidth]{chapter4/images/settings_screen.jpeg}
    \caption{General settings}
  \end{subfigure}
  \hfill
  \begin{subfigure}[b]{0.30\textwidth}
    \includegraphics[width=\textwidth]{chapter4/images/settings_screen_1.jpeg}
    \caption{Security settings}
  \end{subfigure}
  \hfill
  \begin{subfigure}[b]{0.30\textwidth}
    \includegraphics[width=\textwidth]{chapter4/images/settings_screen_2.jpeg}
    \caption{Account management}
  \end{subfigure}

  \caption{Application settings}
  \label{fig:settings}
\end{figure}

\subsection{Real-Time Notifications}
\begin{itemize}
  \item \textbf{Medical file access request:} When a doctor searches for a patient by name and submits an access request to consult the patient's medical file, the patient receives an instant in-app notification banner. The notification displays the doctor's name and speciality, and provides a direct shortcut to the access management screen where the patient can accept or reject the request, as shown in Figures~\ref{fig:notif_access1} and~\ref{fig:notif_access2}.
  \item \textbf{New appointment request:} When a patient books an appointment with a doctor, the doctor receives an instant WebSocket notification displaying the patient's name and the requested date and time Figure~\ref{fig:notif_appt}. This allows the doctor to respond promptly without waiting for a scheduled polling cycle.
\end{itemize}
\begin{figure}[H]
  \centering
  \begin{subfigure}[t]{0.30\textwidth}
    \centering
    \includegraphics[width=\textwidth]{chapter4/images/notif_access_1.jpeg}
    \caption{Doctor requests medical file access (patient view)}
    \label{fig:notif_access1}
  \end{subfigure}
  \hfill
  \begin{subfigure}[t]{0.30\textwidth}
    \centering
    \includegraphics[width=\textwidth]{chapter4/images/notif_access_2.jpeg}
    \caption{Access management screen after receiving the request}
    \label{fig:notif_access2}
  \end{subfigure}
  \hfill
  \begin{subfigure}[t]{0.30\textwidth}
    \centering
    \includegraphics[width=\textwidth]{chapter4/images/notif_appointment.jpeg}
    \caption{New appointment request (doctor view)}
    \label{fig:notif_appt}
  \end{subfigure}
  \caption{Notification and access management screens}
  \label{fig:notif_all}
\end{figure}


\subsection{Male Patient Interface}
Sahhty is not limited to pregnant women. Male patients benefit from a dedicated experience that removes all pregnancy- and menstrual-related content and adopts a distinct blue gradient theme, clearly differentiating the interface from the pink feminine theme used for female patients. The male dashboard focuses on the health indicators relevant to any patient: vital sign measurements, appointment management, alerts, and personal profile.As shown in Figure~\ref{fig:male_interface}, the interface includes five screens: home (Figure~\ref{fig:male_home}), measurements (Figure~\ref{fig:male_measurements}), alerts (Figure~\ref{fig:male_alerts}), appointments (Figure~\ref{fig:male_appointments}), and profile (Figure~\ref{fig:male_profile}).
\begin{figure}[H]
  \centering
  \begin{subfigure}[b]{0.30\textwidth}
    \includegraphics[width=\textwidth]{chapter4/images/home_male.png}
    \caption{Home}
    \label{fig:male_home}
  \end{subfigure}
  \hfill
  \begin{subfigure}[b]{0.30\textwidth}
    \includegraphics[width=\textwidth]{chapter4/images/measurements_male.png}
    \caption{Measurements}
    \label{fig:male_measurements}
  \end{subfigure}
  \hfill
  \begin{subfigure}[b]{0.30\textwidth}
    \includegraphics[width=\textwidth]{chapter4/images/alerts_male.png}
    \caption{Alerts}
    \label{fig:male_alerts}
  \end{subfigure}

  \begin{subfigure}[b]{0.30\textwidth}
    \includegraphics[width=\textwidth]{chapter4/images/appointments_male.png}
    \caption{Appointments}
    \label{fig:male_appointments}
  \end{subfigure}
  \hspace{0.02\textwidth}
  \begin{subfigure}[b]{0.30\textwidth}
    \includegraphics[width=\textwidth]{chapter4/images/profile_male.png}
    \caption{Profile}
    \label{fig:male_profile}
  \end{subfigure}
  \caption{Male patient interface (five main screens)}
  \label{fig:male_interface}
\end{figure}

%\section{AI-Based Pregnancy Risk Assessment}
%In this section, the AI-based pregnancy risk assessment system integrated into the Sahhty platform is presented. The platform uses a trained XGBoost classifier to dynamically evaluate the patient's health condition based on recorded measurements and automatically classify the pregnancy risk level as \textit{low}, \textit{medium}, or \textit{high}.

%\subsection{Risk Assessment Workflow}

%The risk assessment process is automatically triggered whenever a patient submits a new health measurement through the mobile application. The workflow consists of the following steps:

%\begin{enumerate}
 %   \item The patient submits a new measurement from the Flutter application.

 %   \item The measurement is transmitted to the Django backend through the REST API.

  %  \item The backend stores the measurement in the PostgreSQL database and invokes the AI prediction module.

   % \item The prediction engine loads the trained XGBoost model and preprocesses the patient data.

  %  \item The model predicts the pregnancy risk level and returns a probability score.

   % \item The API response includes the generated risk level.

    %\item The Flutter application displays the corresponding visual feedback:
    %\begin{itemize}
     %   \item Orange warning dialog for \textit{medium} risk.
     %   \item Red critical alert dialog for \textit{high} risk.
    %\end{itemize}

    %\item For high-risk situations, the backend additionally generates a persistent alert and sends a push notification through Firebase Cloud Messaging (FCM).
%\end{enumerate}

%\subsection{Technical Configuration}

%The technical characteristics of the AI-based risk assessment engine are summarized in Table~\ref{tab:ai_risk}.

%\begin{table}[H]
%\centering
%\caption{AI-Based Pregnancy Risk Assessment Configuration}
%\label{tab:ai_risk}
%\renewcommand{\arraystretch}{1.3}

%\begin{tabular}{|p{4cm}|p{9cm}|}
%\hline
%\textbf{Component} & \textbf{Description} \\
%\hline
%AI Model & XGBoost Classifier \\
%\hline
%Preprocessing & Scikit-learn preprocessing pipeline \\
%\hline
%Input Features & Weight, blood pressure, blood glucose, heart rate, temperature, and gestational age \\
%\hline
%Prediction Output & Low, Medium, or High pregnancy risk level \\
%\hline
%Execution Trigger & Triggered after each new measurement submission \\
%\hline
%Frontend Feedback & Visual warning and critical alert dialogs \\
%\hline
%Additional Actions & Persistent alert creation and FCM push notification for high-risk cases \\
%\hline
%\end{tabular}
%\end{table}

%--------------------------------------------------
%\section{Smartwatch Integration}
%--------------------------------------------------
%In this section, the Sahhty platform integrates with Wear OS smartwatches through a dedicated companion application to support automatic heart rate monitoring. This integration enables continuous health data collection without requiring direct interaction from the patient.

%\subsection{Integration Architecture}

%The smartwatch integration is based on two main layers:

%\begin{itemize}
 %   \item \textbf{Wear OS Application:} A Kotlin-based smartwatch application responsible for reading heart rate measurements using Android sensor APIs and periodically transmitting the data.

  %  \item \textbf{Flutter Mobile Application:} Receives smartwatch data through the Wearable Data Layer API and forwards the measurements to the Django backend through a Flutter service.
%\end{itemize}

%\subsection{Smartwatch Data Flow}

%The complete smartwatch communication process is summarized in Table~\ref{tab:watch_flow}.

%\begin{table}[H]
%\centering
%\caption{Wear OS Smartwatch Data Transmission Flow}
%\label{tab:watch_flow}
%\renewcommand{\arraystretch}{1.3}

%\begin{tabular}{|p{1cm}|p{4cm}|p{8cm}|}
%\hline
%\textbf{Step} & \textbf{Component} & \textbf{Operation} \\
%\hline
%1 & Wear OS Sensor & Heart rate measurement acquisition \\
%\hline
%2 & WorkManager & Periodic execution every 30 minutes \\
%\hline
%3 & Wearable Data Layer & Transmission of heart rate data to the phone \\
%\hline
%4 & Android Service & Reception of smartwatch data on the mobile device \\
%\hline
%5 & Flutter MethodChannel & Transfer of data to the Flutter layer \\
%\hline
%6 & Flutter Service & Formatting and API request generation \\
%\hline
%7 & Django Backend & Measurement storage and AI risk analysis execution \\
%\hline
%\end{tabular}
%\end{table}

%The smartwatch measurements are transmitted to the backend using a structured JSON payload containing the measurement type, measured value, timestamp context, and patient identifier.

\section{Conclusion}

This chapter presented the implementation of the Sahhty platform, including the adopted technologies, the main user interfaces, and the platform’s innovative features. The developed system provides an intuitive and scalable healthcare solution adapted to the needs of pregnant patients and doctors.

